\section{Directions for Future Research}\label{sec:future}

Our results answer one question---how much of the news is left after the first minute---and
raise several others. We group them by the finding that motivates each, and for each we say what
data or test would answer it.

\paragraph{1. Execution in the first minutes.} Every result in the paper depends on the cost
of a trade one minute after a release, and we model that cost as a fixed number of ticks. In
practice, depth in the order book thins out just before a release and rebuilds over the next
minutes, so slippage is itself state-dependent: it is likely to be largest after the largest
surprises, which are the trades the surprise model takes. Recorded fills, or depth-of-book data
(for example Databento's MBP-10 feed for the same contracts), would allow the cost in
equation~\eqref{eq:cost} to be replaced by a measured cost that depends on the size of the surprise,
the time since the release and the trade's size. The same data would answer how much capital the
strategy can carry before its own trades move the price, a question our backtests cannot address.

\paragraph{2. Price discovery inside the first minute.} We find that most of the response is in
the first minute, but our bars cannot say how it is distributed within that minute. Tick data
would show how many seconds the market needs to price a surprise, whether the speed has changed
over the sample, and whether the residual drift we trade is related to how quickly the first
seconds absorb the news. It would also allow the analysis of latency in
Section~\ref{sec:environment} to be extended from minutes to seconds, so that the value of the
first-minute jump to a faster trader can be measured rather than excluded.

\paragraph{3. Better measures of the surprise.} We measure the surprise as the actual figure minus
the median forecast, standardized. Several refinements are natural. The dispersion of forecasts
measures how uncertain the market was, and a surprise relative to an uncertain consensus may be
priced differently. Revisions to earlier releases and the components of a release (core versus
headline, services versus goods) may carry information that the headline surprise misses. Nowcasts
and market-implied expectations, such as inflation swaps or fed funds futures, are a measure of what
was expected that is less stale than a survey. Finally, the text of the release and of FOMC
statements and minutes can be scored with language models, which would bring the FOMC decision---which
our impact screen does not admit---into the event set.

\paragraph{4. Regimes and the Fed's reaction function.} The surprise book earns most of its return in
2022: without that year its Sharpe ratio is close to zero in every equity market
(Figure~\ref{fig:byyear}). A plausible explanation is that inflation surprises matter most when the
Fed is expected to respond to them. This can be tested directly by conditioning the drift on
measures of the policy regime---the expected path of rates, whether the Fed is tightening, the
level of inflation relative to target---and by letting the coefficients in
equation~\eqref{eq:drift_model} vary over time, for example with a Kalman filter or a rolling
Bayesian prior. A model of this kind would also say in advance when the edge is likely to be small,
which is as valuable to a practitioner as the edge itself.

\paragraph{5. Extensions of price-path analogues.} DTW analogues do not forecast the direction of
the drift, but their dispersion forecasts its size, the first window after the release is the best of the five, and the analogues add value when they agree with the surprise:
in \YM\ the book that trades only when the two agree has a Sharpe ratio of 0.94 with the same value
outside 2022. These findings suggest four extensions. First, multivariate paths that match the equity
and Treasury reaction together, since the joint reaction says more about the nature of the news
than either market alone. Second, a book that uses DTW only for sizing and the surprise for direction,
tested out of sample on all four markets. Third, learned similarity measures (for example shapelets or
a neural embedding of the path) in place of a fixed distance. Fourth, a book that trades only the
first window, which keeps the part of the ladder with information and drops the part that only adds
costs.

\paragraph{6. Other markets and countries.} The tick-size explanation for the Treasury result
predicts that a contract with a finer tick relative to its moves should show a tradable drift. This
can be tested on the 2-year and 5-year notes and on SOFR futures, whose ticks are small relative to
their moves around policy news, and on the Bund, the gilt and JGB futures around their own countries'
releases. In equities, the Russell 2000 future (traded since 2017) and the Euro Stoxx 50 offer a
test of whether the drift is larger in less liquid or more rate-sensitive indices. Currency and
commodity futures around U.S.\ releases would show whether the drift is a feature of U.S.\ equity
futures or of how markets in general absorb scheduled news.

\paragraph{7. Portfolio construction.} Our portfolios use fixed equal-risk weights, chosen on
2016--2018. The approaches lose in different years, which suggests that weights that adapt---to recent
performance, to the regime of item 4, or to the expected cost of item 1---could do better, though at
the price of more choices and a higher risk of overfitting. A careful test would fix the rule for
changing the weights in advance and report it with the same random-sign and walk-forward checks we
use for the books themselves.

\paragraph{8. Why does the drift exist?} We document a drift but test its cause only indirectly.
Three explanations make different predictions. Limits to arbitrage predict a larger drift when
liquidity is thin and costs are high. Limited attention predicts a larger drift when several releases
arrive at the same time or when the surprise is in a less-watched component. Disagreement predicts a
larger drift when forecasts are dispersed and when volume after the release is high. Volume, order-flow
imbalance and the dispersion of forecasts around each release are enough to tell them apart, and the
answer would say whether the edge is likely to last.
